\section{Definitions and Assumptions}

\begin{definition}[Plan]
A \textbf{plan} for transition $s_0 \to s_1$ is a pre-action specification of
at least one future decision window, one entry action, and one relevant
environmental arrangement for the target state.
\end{definition}

\begin{definition}[Resolved entry specification]
Let $U(s_1)$ denote the target ambiguity from
Chapter~\ref{ch:formal-language}. A plan is said to provide a
\textbf{resolved entry specification} if it fixes a visible first action for
entering $s_1$ and thereby lowers the unresolved specification burden.
\end{definition}

\begin{definition}[Execution choice complexity]
For a given future decision window, let $D$ denote the number of
action-relevant micro-decisions still left to be made at the moment of action.
Examples include choosing the location, selecting the first task, locating the
materials, or deciding when to stop.
\end{definition}

\begin{remark}
The variable $D$ is introduced as a qualitative bookkeeping device, not a
measured cognitive constant. The only property needed below is monotonicity:
if more in-the-moment decisions remain unresolved, then the execution burden is
harder rather than easier.
\end{remark}

\begin{axiom}[Monotone planning assumption]
Holding the current state fixed, the activation barrier
$\DV(s_0 \to s_1)$ is weakly increasing in target ambiguity and weakly
decreasing in relevant preparation.
\end{axiom}

\begin{axiom}[Decision-load assumption]
Holding other terms fixed, the in-the-moment execution burden is weakly
increasing in $D$.
\end{axiom}

\begin{discussion}
These assumptions do not prove that a specific numerical planner outperforms a
specific spontaneous agent in every situation. They only formalize the business
logic already used by the book: leaving fewer questions unanswered makes future
entry easier, and resolving relevant setup in advance lowers the burden at the
moment of action.
\end{discussion}
